% GENERATED FILE. Edit canonical lessons, then run npm run generate:books.
% canonical-source-hash: fnv1a64:ebf65cdefdf82b6e
% canonical-lessons: MR-C305-gatvarsi, MR-C305-magil, MR-C305-tatpurvi, MR-C305-kalantarane, MR-C305-lotne

\chapter{Last Year and Before}
\label{ch:mr-last-year-and-before}
\hlchaptermodality{305}

\begin{chapteropening}
\textbf{By the end of this chapter:} \emph{I can place events in the past: last year, before that, in time.}

Complete the last lesson of chapter 305: I can place events in the past: last year, before that, in time.
\end{chapteropening}

\section[\texorpdfstring{\mr{गतवर्षी} (gatvarṣī)}{गतवर्षी (gatvarṣī)}]{\mr{गतवर्षी} (gatvarṣī) — last year}

\label{lesson:MR-C305-gatvarsi}



\begin{quote}
Before the new one: say the Marathi for to make a stir, then the Marathi for to spread.
\end{quote}

\subsection*{You'll want to know: \mr{गतवर्षी}}
\textbf{\mr{गतवर्षी}} — \emph{gatvarṣī} — \textquotedblleft{}last year\textquotedblright{}.

\begin{quote}
\mr{गतवर्षी} \mr{आमच्या} \mr{गावात} \mr{मोठी} \mr{जत्रा} \mr{भरली} \mr{होती}.
\end{quote}

\begin{quote}
\emph{gatvarṣī āmacyā gāvāt moṭhī jatrā bharlī hotī.}
\end{quote}

\begin{quote}
Last year a big fair was held in our village.
\end{quote}

Formal; in speech \mr{मागच्या} \mr{वर्षी}. The letter \mr{ञ} hides inside \mr{ज्ञ}: \mr{ज्ञान}, knowledge.

\begin{cousinweb}
\textbf{\mr{गतवर्षी}} (\emph{gatvarṣī}) is built from Sanskrit \emph{gata}, 'gone', and \emph{varṣa}, 'a year' (Monier-Williams), in the Marathi locative: in the year gone by.
\end{cousinweb}

\begin{grammarlens}[title={What you've built}]
One more word for telling what happened.
\end{grammarlens}

\subsection*{Your turn}
\begin{itemize}
\raggedright
  \item \emph{Say it:} \emph{gatvarṣī}
  \item \emph{Say it:} \emph{gatvarṣī}, once more
  \item \emph{Say it:} \emph{gatvarṣī āmacyā gāvāt moṭhī jatrā bharlī hotī.} — Last year a big fair was held in our village.
  \item \emph{Recall:} say the Marathi for to make a stir, then the Marathi for to spread, then say \emph{gatvarṣī} again
\end{itemize}

\begin{tcolorbox}[breakable,colback=teal!4,colframe=teal!35!black,title={Before you move on}]
What does \mr{गतवर्षी} mean? (\textquotedblleft{}Last year\textquotedblright{}.) Say it once more.
\end{tcolorbox}
\section[\texorpdfstring{\mr{मागील} (māgīl)}{मागील (māgīl)}]{\mr{मागील} (māgīl) — previous; the last}

\label{lesson:MR-C305-magil}



\begin{quote}
Before the new one: say the Marathi for to spread, then the Marathi for last year.
\end{quote}

\subsection*{You'll want to know: \mr{मागील}}
\textbf{\mr{मागील}} — \emph{māgīl} — \textquotedblleft{}previous; the last\textquotedblright{}.

\begin{quote}
\mr{मागील} \mr{बैठकीत} \mr{हा} \mr{प्रश्न} \mr{आला} \mr{होता}.
\end{quote}

\begin{quote}
\emph{māgīl baiṭhkīt hā praśna ālā hotā.}
\end{quote}

\begin{quote}
This question came up at the previous meeting.
\end{quote}

Formal; it keeps one form: \mr{मागील} \mr{आठवडा}, \mr{मागील} \mr{बैठक}.

\begin{cousinweb}
\textbf{\mr{मागील}} (\emph{māgīl}) is built on \emph{māg}, 'a track; behind', which CDIAL lists as Marathi 'road, track' under Sanskrit \emph{mārga} (CDIAL 10071), with the ending \emph{-īl}.
\end{cousinweb}

\begin{grammarlens}[title={What you've built}]
One more word for telling what happened.
\end{grammarlens}

\subsection*{Your turn}
\begin{itemize}
\raggedright
  \item \emph{Say it:} \emph{māgīl}
  \item \emph{Say it:} \emph{māgīl}, once more
  \item \emph{Say it:} \emph{māgīl baiṭhkīt hā praśna ālā hotā.} — This question came up at the previous meeting.
  \item \emph{Recall:} say the Marathi for to spread, then the Marathi for last year, then say \emph{māgīl} again
\end{itemize}

\begin{tcolorbox}[breakable,colback=teal!4,colframe=teal!35!black,title={Before you move on}]
What does \mr{मागील} mean? (\textquotedblleft{}Previous; the last\textquotedblright{}.) Say it once more.
\end{tcolorbox}
\section[\texorpdfstring{\mr{तत्पूर्वी} (tatpūrvī)}{तत्पूर्वी (tatpūrvī)}]{\mr{तत्पूर्वी} (tatpūrvī) — before that}

\label{lesson:MR-C305-tatpurvi}



\begin{quote}
Before the new one: say the Marathi for last year, then the Marathi for previous.
\end{quote}

\subsection*{You'll want to know: \mr{तत्पूर्वी}}
\textbf{\mr{तत्पूर्वी}} — \emph{tatpūrvī} — \textquotedblleft{}before that\textquotedblright{}.

\begin{quote}
\mr{कार्यक्रम} \mr{सहा} \mr{वाजता} \mr{होता}, \mr{तत्पूर्वी} \mr{सगळे} \mr{जेवले}.
\end{quote}

\begin{quote}
\emph{kāryakram sahā vājtā hotā, tatpūrvī sagḷe jevle.}
\end{quote}

\begin{quote}
The programme was at six; before that everyone ate.
\end{quote}

Formal and written; in speech \mr{त्याच्या} \mr{आधी}.

\begin{cousinweb}
\textbf{\mr{तत्पूर्वी}} (\emph{tatpūrvī}) is built from Sanskrit \emph{tat}, 'that', and \emph{pūrva}, 'before' (Monier-Williams), with a Marathi locative ending.
\end{cousinweb}

\begin{grammarlens}[title={What you've built}]
One more word for telling what happened.
\end{grammarlens}

\subsection*{Your turn}
\begin{itemize}
\raggedright
  \item \emph{Say it:} \emph{tatpūrvī}
  \item \emph{Say it:} \emph{tatpūrvī}, once more
  \item \emph{Say it:} \emph{kāryakram sahā vājtā hotā, tatpūrvī sagḷe jevle.} — The programme was at six; before that everyone ate.
  \item \emph{Recall:} say the Marathi for last year, then the Marathi for previous, then say \emph{tatpūrvī} again
\end{itemize}

\begin{tcolorbox}[breakable,colback=teal!4,colframe=teal!35!black,title={Before you move on}]
What does \mr{तत्पूर्वी} mean? (\textquotedblleft{}Before that\textquotedblright{}.) Say it once more.
\end{tcolorbox}
\section[\texorpdfstring{\mr{कालांतराने} (kālāntarāne)}{कालांतराने (kālāntarāne)}]{\mr{कालांतराने} (kālāntarāne) — over time; as time went by}

\label{lesson:MR-C305-kalantarane}



\begin{quote}
Before the new one: say the Marathi for previous, then the Marathi for before that.
\end{quote}

\subsection*{You'll want to know: \mr{कालांतराने}}
\textbf{\mr{कालांतराने}} — \emph{kālāntarāne} — \textquotedblleft{}over time; as time went by\textquotedblright{}.

\begin{quote}
\mr{कालांतराने} \mr{गावाचं} \mr{शहर} \mr{झालं}.
\end{quote}

\begin{quote}
\emph{kālāntarāne gāvācaṁ śahar jhālaṁ.}
\end{quote}

\begin{quote}
As time went by, the village became a town.
\end{quote}

Built from Sanskrit kāla, time, and antara, an interval, with -\mr{ने}: as time went by.

\begin{cousinweb}
\textbf{\mr{कालांतराने}} (\emph{kālāntarāne}) is built on Sanskrit \emph{kālāntara}, 'an interval of time' (Monier-Williams), with the Marathi instrumental \emph{-ne}.
\end{cousinweb}

\begin{grammarlens}[title={What you've built}]
One more word for telling what happened.
\end{grammarlens}

\subsection*{Your turn}
\begin{itemize}
\raggedright
  \item \emph{Say it:} \emph{kālāntarāne}
  \item \emph{Say it:} \emph{kālāntarāne}, once more
  \item \emph{Say it:} \emph{kālāntarāne gāvācaṁ śahar jhālaṁ.} — As time went by, the village became a town.
  \item \emph{Recall:} say the Marathi for previous, then the Marathi for before that, then say \emph{kālāntarāne} again
\end{itemize}

\begin{tcolorbox}[breakable,colback=teal!4,colframe=teal!35!black,title={Before you move on}]
What does \mr{कालांतराने} mean? (\textquotedblleft{}Over time; as time went by\textquotedblright{}.) Say it once more.
\end{tcolorbox}
\section[\texorpdfstring{\mr{लोटणे} (loṭṇe)}{लोटणे (loṭṇe)}]{\mr{लोटणे} (loṭṇe) — to elapse; to push}

\label{lesson:MR-C305-lotne}



\begin{quote}
Before the new one: say the Marathi for before that, then the Marathi for over time.
\end{quote}

\subsection*{You'll want to know: \mr{लोटणे}}
\textbf{\mr{लोटणे}} — \emph{loṭṇe} — \textquotedblleft{}to elapse; to push\textquotedblright{}.

\begin{quote}
\mr{त्या} \mr{घटनेला} \mr{आता} \mr{दहा} \mr{वर्षं} \mr{लोटली}.
\end{quote}

\begin{quote}
\emph{tyā ghaṭnelā ātā dahā varṣaṁ loṭlī.}
\end{quote}

\begin{quote}
Ten years have now passed since that event.
\end{quote}

Past (he …): \mr{तो} \mr{लोटला}, \emph{to loṭlā}.

Of time, with -\mr{ला}: \mr{घटनेला} \mr{वर्षं} \mr{लोटली}, years have passed since the event. Also to push: \mr{दार} \mr{लोटणे}.

\begin{cousinweb}
\textbf{\mr{लोटणे}} (\emph{loṭṇe}) is inherited: CDIAL lists Marathi \emph{loṭṇẽ}, 'to roll, wallow', under \emph{*lōrtati}, 'rolls' (CDIAL 11156). The senses 'to push' and, of time, 'to pass' are most likely the same verb: time rolls on.
\end{cousinweb}

\begin{grammarlens}[title={What you've built}]
One more word for telling what happened.
\end{grammarlens}

\subsection*{Your turn}
\begin{itemize}
\raggedright
  \item \emph{Say it:} \emph{loṭṇe}
  \item \emph{Say it:} \emph{loṭṇe}, once more
  \item \emph{Say it:} \emph{tyā ghaṭnelā ātā dahā varṣaṁ loṭlī.} — Ten years have now passed since that event.
  \item \emph{Recall:} say the Marathi for before that, then the Marathi for over time, then say \emph{loṭṇe} again
\end{itemize}

\begin{tcolorbox}[breakable,colback=teal!4,colframe=teal!35!black,title={Before you move on}]
What does \mr{लोटणे} mean? (\textquotedblleft{}To elapse; to push\textquotedblright{}.) Say it once more.
\end{tcolorbox}
